
@@M0@@
l'application qui associe à la classe @@M1@@ d'un @@M2@@-module inversible @@M3@@ l'élément @@M4@@ de @@M5@@. Cette application est évidemment un homomorphisme de @@M6@@ dans le groupe multiplicatif @@M7@@ des éléments inversibles de @@M8@@. Par la suite nous noterons multiplicativement la loi de groupe de @@M9@@.

L'application @@M10@@ a un inverse à gauche défini par le foncteur «~det~» : Lorsque @@M11@@ est un @@M12@@-module localement libre de rang @@M13@@, on a
@@M14@@
Si
@@M15@@
est une suite exacte de @@M16@@-modules localement libres, il y a un isomorphisme canonique
@@M17@@
provenant de (6.2.9) et du fait que (5.1.2) est localement scindée. On obtient donc un homomorphisme de groupes abéliens
@@M18@@
qui est l'inverse à gauche cherché de @@M19@@. En particulier @@M20@@ est injective, et nous identifierons par la suite @@M21@@ avec son image dans @@M22@@.

\subsection*{5.2.}
On définit une autre application
@@M23@@
